\subsection{Temporal Video Grounding}

Video Temporal Grounding (VTG) aims to localize the temporal segment in a video that corresponds to a language query. Recent progress in VTG increasingly leverages multimodal large language models (MLLMs), shifting from task-specific localization pipelines toward more unified video-language frameworks \cite{wu2025survey_vtg_mllm}. Representative methods such as VTimeLLM \cite{huang2024vtimellm}, TimeChat \cite{ren2024timechat}, and LITA \cite{huang2024lita} improve timestamp sensitivity, temporal representation, and localization-oriented instruction tuning for video understanding. Grounded-VideoLLM further strengthens fine-grained localization through grounding-oriented architectural design and multi-stage training \cite{wang2025grounded_videollm}. 

Beyond supervised grounding, training-free and post-training paradigms have also shown promising results. Training-free VTG methods demonstrate that large pre-trained models already contain useful temporal reasoning priors \cite{zheng2024training,xu2025zero,qu2024chatvtg}, while more recent approaches such as Time-R1 and TimeLens highlight the importance of grounding-specific post-training, stronger data curation, and more reliable optimization for precise temporal localization \cite{wang2025timer1,zhang2026timelens}. Overall, these works substantially improve the ability of MLLMs to localize relevant temporal moments. However, they mainly study VTG in the conventional query-to-moment setting. In contrast, our work focuses on \emph{multiple-choice grounded VideoQA}, where the input answer candidates are not naturally written as retrieval-friendly grounding queries, creating an additional mismatch between answer formulation and temporal evidence acquisition.

\subsection{Grounded and Reasoning VideoQA}

Conventional VideoQA primarily evaluates answer correctness, often without requiring the model to identify the evidence supporting its prediction. This limitation has motivated grounded VideoQA, where correct answering must be accompanied by faithful temporal evidence \cite{xiao2024can}. NExT-GQA is a representative benchmark in this direction, showing that strong answer accuracy does not necessarily imply grounded reasoning \cite{xiao2024can}. ReXTime complements this setting by emphasizing reasoning across time, where the evidence needed for answering may appear in a different temporal region from the question context \cite{chen2024rextime}. More broadly, recent empirical studies in the LLM era also suggest that answer accuracy alone is insufficient for assessing genuine video understanding \cite{xiao2025videoqa}.

Several recent works attempt to connect temporal localization and answer prediction more explicitly. SeViLA decomposes the task into localization and answering stages, reducing reliance on dense temporal annotations through self-refinement \cite{yu2023sevila}. For long-video reasoning, LLoVi, LangRepo, and VideoStreaming show that staged language-centric processing, memory-like representations, and streaming mechanisms can improve reasoning over temporally extended videos \cite{zhang2024simple,kahatapitiya2025language,qian2024streaming}. More recently, VideoExpert and VideoMind demonstrate that stronger temporal grounding and role specialization can further benefit grounded reasoning, especially in long-form and temporally complex settings \cite{zhao2025videoexpert,liu2025videomind}. 

Our work is most closely related to this line of grounded reasoning methods, but differs in its task-specific focus. Existing approaches typically connect localization and answering more directly, whereas we target a distinct bottleneck in \emph{multiple-choice} grounded VideoQA: answer options are designed for discrimination rather than temporal evidence retrieval. To address this mismatch, we explicitly introduce a refinement stage to transform answer candidates into retrieval-friendly queries and a validation stage to score and rank retrieved evidence before final answer generation.
